\subsection{一次变速先确定新轨道，再讨论怎样接续}
瞬时变速只改变速度，不改变该瞬间的位置。此后若发动机关闭，卫星沿由新能量、新角动量确定的轨道运动；不能把“速度变大”直接翻译成“立即进入更大的圆”。较高圆轨道的圆周速度虽然更低，从较低圆轨道转移过去却通常要先顺行加速，因为首先要提高轨道能量。

\begin{origexamplebox}{两个共面圆轨道之间的双脉冲转移}\label{ch11:ex:teaching-transfer}
卫星原在半径 $r_1$ 的圆轨道上，拟转入同平面、同绕行方向、半径 $r_2>r_1$ 的圆轨道。采用一条近点为 $r_1$、远点为 $r_2$ 的椭圆转移轨道，在两端各施加一次沿切向的瞬时速度增量。忽略卫星质量及推进过程的有限持续时间，求两次增量与转移时间。
\end{origexamplebox}
\begin{xesolution}
图~\ref{ch11:fig:apsis} 的 $P,A$ 分别对应转移椭圆的近点、远点，故半长轴 $a=(r_1+r_2)/2$。记 $\mu=GM$。转移过程中单位质量能量为 $-\mu/(2a)$，在任一点有
\[
 \frac{v^2}{2}-\frac\mu r=-\frac\mu{2a},
 \qquad v^2=\mu\left(\frac2r-\frac1a\right).
\]
这先由能量守恒给出端点速度，尚未把它们误认为圆周速度。在 $P$ 点：
\[
 v_P=\sqrt{\frac{2\mu r_2}{r_1(r_1+r_2)}},\qquad
 \Delta v_1=v_P-\sqrt{\frac\mu{r_1}}>0.
\]
第一次顺行加速，把圆轨道切换为转移椭圆。飞到 $A$ 点时：
\[
 v_A=\sqrt{\frac{2\mu r_1}{r_2(r_1+r_2)}},\qquad
 \Delta v_2=\sqrt{\frac\mu{r_2}}-v_A>0.
\]
第二次仍是顺行加速，因为转移椭圆远点速度比当地圆周速度更小。两端满足 $r_1v_P=r_2v_A$，核验了无推力飞行段的角动量守恒；两次脉冲处卫星自身角动量则改变，不能跨过脉冲继续当作守恒。
从近点到远点用半个椭圆周期，因此
\[
 t_{\rm tr}=\pi\sqrt{\frac{a^3}{\mu}}
 =\pi\sqrt{\frac{(r_1+r_2)^3}{8\mu}}.
\]
$r_2\to r_1$ 时两次速度增量均趋零；形式上的飞行时间趋于原圆轨道半周期，表示仍按“飞到对面”的约定计时，而不是说留在原轨道必须等待半周期。结果的速度与时间量纲分别由 $\sqrt{\mu/r}$、$\sqrt{r^3/\mu}$ 给出。
\end{xesolution}
